%========================================
% MatinBook - Stage 14: Complete Book Test
% Test: ALL Modules + Bibliography + Index + Cover
% Purpose: Full integration test of all MatinBook features
%========================================

\makeatletter
\def\input@path{{../}}
\makeatother

\documentclass{matinbook}

%----------------------------------------
% Bibliography
%----------------------------------------
\addbibresource{../references.bib}

%----------------------------------------
% Book Metadata
%----------------------------------------
\title{راهنمای جامع برنامه‌نویسی و الگوریتم}
\author{تیم توسعه MatinBook}
\date{تابستان ۱۴۰۵}

\begin{document}

%========================================
% COVER
%========================================
\makecover
{راهنمای جامع\\برنامه‌نویسی و الگوریتم}
{از مفاهیم پایه تا پیشرفته}
{تیم توسعه MatinBook}
{تابستان ۱۴۰۵}

\makebackcover{%
    این کتاب راهنمای کاملی برای یادگیری برنامه‌نویسی
    و الگوریتم است که با استفاده از کلاس \lr{MatinBook}
    نسخه ۱ حروف‌چینی شده است.

    \textbf{ویژگی‌های کتاب:}
    \begin{itemize}
        \item آموزش گام‌به‌گام برنامه‌نویسی پایتون
        \item پیاده‌سازی الگوریتم‌های کلاسیک
        \item تحلیل پیچیدگی و مقایسه الگوریتم‌ها
        \item نمودارها و شبه‌کدهای استاندارد
        \item تمرین‌های متنوع با راه حل تشریحی
        \item منابع معتبر و نمایه جامع
    \end{itemize}

    مناسب برای دانشجویان، مدرسان و علاقه‌مندان
    به علوم کامپیوتر.

    \textbf{ساخته شده با \lr{MatinBook} نسخه ۱}
}

%========================================
% TITLE PAGE
%========================================
\maketitle

%========================================
% COPYRIGHT
%========================================
\thispagestyle{empty}
\vspace*{6cm}
\begin{center}
    {\large\textbf{حقوق نشر}}\\[1cm]
    تمامی حقوق این کتاب برای تیم \lr{MatinBook} محفوظ است.\\
    هرگونه کپی‌برداری، تکثیر و انتشار بدون اجازه کتبی ممنوع است.\\[1cm]
    نسخه ۱.۰ --- تابستان ۱۴۰۵\\[1cm]
    \lr{Copyright \copyright\ 2026 MatinBook Project}\\
    \lr{Licensed under the MIT License}
\end{center}
\clearpage

%========================================
% DEDICATION
%========================================
\thispagestyle{empty}
\vfill
\begin{flushright}
    {\Large\itshape
        تقدیم به تمام جویندگان دانش\\
        که هیچ‌گاه از یادگیری\\
        دست نمی‌کشند.
    }
\end{flushright}
\clearpage
\frontmatter

%========================================
% PREFACE
%========================================
\chapter*{پیشگفتار}
\addcontentsline{toc}{chapter}{پیشگفتار}

\section*{درباره این کتاب}

این کتاب حاصل ماه‌ها تلاش تیم \lr{MatinBook}
برای ایجاد یک منبع آموزشی کامل در زمینه
برنامه‌نویسی و الگوریتم است. هدف ما ارائه
مطالبی دقیق، روان و کاربردی با استفاده از
ابزارهای حروف‌چینی مدرن فارسی است.

\section*{ساختار کتاب}

کتاب در سه بخش اصلی سازماندهی شده است:

\begin{enumerate}
    \item \textbf{بخش اول --- مبانی برنامه‌نویسی:}
    آشنایی با پایتون، ساختارهای داده و کنترل جریان
    
    \item \textbf{بخش دوم --- الگوریتم‌ها:}
    الگوریتم‌های جستجو، مرتب‌سازی و تحلیل پیچیدگی
    
    \item \textbf{بخش سوم --- مفاهیم پیشرفته:}
    برنامه‌نویسی شیءگرا، ساختمان داده‌های پیشرفته
    و الگوریتم‌های گراف
\end{enumerate}

\section*{نحوه استفاده}

هر فصل شامل توضیحات نظری، کدهای نمونه،
نمودارهای آموزشی و تمرین‌های عملی است.
توصیه می‌شود خواننده علاوه بر مطالعه متن،
کدها را اجرا و تمرین‌ها را حل کند.

\vspace{1cm}
\begin{flushleft}
    \textbf{تیم توسعه MatinBook}\\
    تابستان ۱۴۰۵
\end{flushleft}

\clearpage

%========================================
% TABLE OF CONTENTS
%========================================
\tableofcontents
\clearpage

\listoffigures
\clearpage

\listoftables
\clearpage

\listofalgorithms
\clearpage
\mainmatter

%========================================
% PART 1: PROGRAMMING BASICS
%========================================
\part{مبانی برنامه‌نویسی}

%========================================
% CHAPTER 1: Python Fundamentals
%========================================
\chapter{مبانی پایتون}
\label{chap:python}

\section{معرفی زبان پایتون}

پایتون\index{پایتون} یک زبان برنامه‌نویسی
سطح بالا، تفسیری و همه‌منظوره است که توسط
خیدو فان روسوم\index{فان روسوم، خیدو}
در سال ۱۹۹۱ میلادی ساخته شد. این زبان
بر خوانایی کد و سادگی نحو تأکید دارد
و امروزه در حوزه‌های مختلفی از جمله
علم داده\index{علم داده}،
هوش مصنوعی\index{هوش مصنوعی}،
توسعه وب\index{توسعه وب} و
خودکارسازی\index{خودکارسازی} استفاده می‌شود.

\begin{definition}[پایتون]
    پایتون یک زبان برنامه‌نویسی شیءگرا،
    تفسیری و با نوع‌دهی پویا (\lr{dynamic typing})
    است که فلسفه طراحی آن بر خوانایی کد
    و کاهش پیچیدگی نحو استوار است.
    \label{def:python}
\end{definition}

\section{انواع داده و متغیرها}

\begin{latin}
\begin{minted}{python}
# Basic data types in Python
name: str = "Ali"                    # String
age: int = 25                        # Integer
height: float = 1.75                 # Float
is_student: bool = True              # Boolean

# Collection types
grades: list = [18, 17, 20, 19]     # List
point: tuple = (10, 20)              # Tuple
info: dict = {"name": "Ali"}         # Dictionary
tags: set = {"python", "coding"}     # Set

# Print all values
print(f"Name: {name}, Age: {age}")
print(f"Height: {height}, Student: {is_student}")
\end{minted}
\end{latin}
\codecaption{انواع داده پایه در پایتون}
\label{code:data-types}

\section{ساختارهای کنترلی}

\begin{definition}[ساختار شرطی]
    ساختار شرطی\index{ساختار شرطی} به برنامه
    اجازه می‌دهد بر اساس شرایط مختلف،
    مسیرهای متفاوتی را اجرا کند.
    \label{def:conditional}
\end{definition}

\begin{latin}
\begin{minted}{python}
def classify_number(n: int) -> str:
    """
    Classify a number as positive, negative, or zero.
    
    Args:
        n: An integer to classify.
    
    Returns:
        A string describing the number.
    """
    if n > 0:
        return "Positive"
    elif n < 0:
        return "Negative"
    else:
        return "Zero"

# Test with various inputs
test_numbers = [5, -3, 0, 42, -17]
for num in test_numbers:
    result = classify_number(num)
    print(f"{num:4d} -> {result}")
\end{minted}
\end{latin}
\codecaption{ساختار شرطی در پایتون}
\label{code:conditional}

\section{حلقه‌ها و تکرار}

\begin{latin}
\begin{minted}{python}
# For loop with range
print("Counting with for loop:")
for i in range(5):
    print(f"  Iteration {i}")

# While loop
print("\nCounting with while loop:")
count = 0
while count < 5:
    print(f"  Count: {count}")
    count += 1

# List comprehension
squares = [x**2 for x in range(10)]
print(f"\nSquares: {squares}")
\end{minted}
\end{latin}
\codecaption{انواع حلقه‌ها در پایتون}
\label{code:loops}

%========================================
% CHAPTER 2: Functions and Modules
%========================================
\chapter{توابع و ماژول‌ها}
\label{chap:functions}

\section{تعریف تابع}

تابع\index{تابع} بلوکی از کد است که
می‌تواند ورودی دریافت کند، عملیاتی
روی آن انجام دهد و خروجی تولید کند.

\begin{definition}[تابع در پایتون]
    تابع با کلمه کلیدی \inlcode{def} تعریف می‌شود
    و می‌تواند پارامترهای ورودی، مستندات
    (\lr{docstring}) و مقدار بازگشتی داشته باشد.
    \label{def:function}
\end{definition}

\begin{latin}
\begin{minted}{python}
def factorial(n: int) -> int:
    """
    Calculate the factorial of n recursively.
    
    Args:
        n: A non-negative integer.
    
    Returns:
        n! = n * (n-1) * ... * 1
    
    Raises:
        ValueError: If n is negative.
    """
    if n < 0:
        raise ValueError("Factorial is not defined for negative numbers")
    if n <= 1:
        return 1
    return n * factorial(n - 1)


# Calculate and display factorials
for i in range(10):
    print(f"{i}! = {factorial(i)}")
\end{minted}
\end{latin}
\codecaption{تابع فاکتوریل بازگشتی}
\label{code:factorial}

\begin{example}
    خروجی تابع \coderef{code:factorial}:
    \begin{align*}
        0! &= 1 \\
        5! &= 5 \times 4 \times 3 \times 2 \times 1 = 120 \\
        10! &= 3628800
    \end{align*}
    \label{ex:factorial-values}
\end{example}

\section{آرگومان‌های تابع}

\begin{latin}
\begin{minted}{python}
def describe_person(
    name: str,
    age: int = 0,
    *hobbies: str,
    **details
) -> str:
    """
    Create a description of a person.
    
    Args:
        name: Person's full name.
        age: Person's age (default: 0).
        *hobbies: Variable number of hobbies.
        **details: Additional key-value details.
    
    Returns:
        A formatted description string.
    """
    parts = [f"Name: {name}"]
    
    if age > 0:
        parts.append(f"Age: {age}")
    
    if hobbies:
        parts.append(f"Hobbies: {', '.join(hobbies)}")
    
    for key, value in details.items():
        parts.append(f"{key}: {value}")
    
    return "\n".join(parts)


# Test the function
print(describe_person("Ali", 25, "coding", "reading"))
print("---")
print(describe_person("Sara", job="Engineer", city="Tehran"))
\end{minted}
\end{latin}
\codecaption{آرگومان‌های متغیر در پایتون}
\label{code:arguments}

%========================================
% PART 2: ALGORITHMS
%========================================
\part{الگوریتم‌ها}

%========================================
% CHAPTER 3: Search Algorithms
%========================================
\chapter{الگوریتم‌های جستجو}
\label{chap:search}

\section{جستجوی خطی}

الگوریتم جستجوی خطی\index{الگوریتم!جستجوی خطی}
ساده‌ترین روش جستجو است. در این روش،
تک‌تک عناصر آرایه بررسی می‌شوند
تا عنصر مورد نظر پیدا شود.

\begin{latin}
\begin{algorithm}
\begin{algorithmic}[1]
    \Require $A[1..n]$: array of $n$ elements
    \Require $x$: target value to find
    \Ensure Index of $x$ in $A$, or $-1$ if not found
    
    \For{$i \gets 1$ \textbf{to} $n$}
        \If{$A[i] = x$}
            \State \Return $i$
        \EndIf
    \EndFor
    \State \Return $-1$
\end{algorithmic}
\end{algorithm}
\algcaption[alg:linear]{جستجوی خطی}
\end{latin}

\begin{theorem}[پیچیدگی جستجوی خطی]
    الگوریتم \algref{alg:linear} در بدترین حالت
    $O(n)$ مقایسه انجام می‌دهد.
    \label{thm:linear-complexity}
\end{theorem}

\section{جستجوی دودویی}

برای آرایه‌های \textbf{مرتب}،
جستجوی دودویی\index{الگوریتم!جستجوی دودویی}
بسیار کارآمدتر است.

\begin{latin}
\begin{algorithm}
\begin{algorithmic}[1]
    \Require $A[1..n]$: sorted array (ascending)
    \Require $x$: target value to find
    \Ensure Index of $x$ in $A$, or $-1$ if not found
    
    \State $left \gets 1$, $right \gets n$
    \While{$left \leq right$}
        \State $mid \gets \lfloor (left + right) / 2 \rfloor$
        \If{$A[mid] = x$}
            \State \Return $mid$
        \ElsIf{$A[mid] < x$}
            \State $left \gets mid + 1$
        \Else
            \State $right \gets mid - 1$
        \EndIf
    \EndWhile
    \State \Return $-1$
\end{algorithmic}
\end{algorithm}
\algcaption[alg:binary]{جستجوی دودویی}
\end{latin}

\begin{figure}[h]
\centering
\begin{latin}
\begin{tikzpicture}
  \begin{axis}[
    width=12cm, height=7cm,
    xlabel={Input Size ($n$)},
    ylabel={Number of Comparisons},
    title={Search Algorithm Comparison},
    grid=major,
    legend pos=north west,
    domain=0:100,
    samples=100
  ]
    \addplot[matinred, thick] {x};
    \addlegendentry{Linear Search $O(n)$}
    \addplot[matinblue, thick] {log2(x)};
    \addlegendentry{Binary Search $O(\log n)$}
  \end{axis}
\end{tikzpicture}
\end{latin}
\caption{مقایسه رشد پیچیدگی الگوریتم‌های جستجو}
\label{fig:search-comparison}
\end{figure}

\begin{remark}
    جستجوی دودویی (\algref{alg:binary})
    فقط روی آرایه‌های \textbf{مرتب} کار می‌کند.
    اگر آرایه مرتب نباشد، ابتدا باید آن را
    مرتب کرد که خود هزینه $O(n \log n)$ دارد.
    \label{rem:binary-requirement}
\end{remark}

%========================================
% CHAPTER 4: Sorting Algorithms
%========================================
\chapter{الگوریتم‌های مرتب‌سازی}
\label{chap:sorting}

\section{مرتب‌سازی سریع}

\lr{Quick Sort}\index{الگوریتم!مرتب‌سازی سریع}
از الگوریتم‌های تقسیم و حل\index{تقسیم و حل}
است که در عمل بسیار سریع عمل می‌کند.

\begin{latin}
\begin{algorithm}
\begin{algorithmic}[1]
    \Function{QuickSort}{$A$, $low$, $high$}
        \If{$low < high$}
            \State $p \gets \Call{Partition}{A, low, high}$
            \State $\Call{QuickSort}{A, low, p-1}$
            \State $\Call{QuickSort}{A, p+1, high}$
        \EndIf
    \EndFunction
\end{algorithmic}
\end{algorithm}
\algcaption[alg:quicksort]{مرتب‌سازی سریع}
\end{latin}

\section{الگوریتم اقلیدس}

\begin{latin}
\begin{algorithm}
\begin{algorithmic}[1]
    \Require Non-negative integers $a, b$
    \Ensure $\gcd(a, b)$
    \While{$b \neq 0$}
        \State $r \gets a \bmod b$
        \State $a \gets b$
        \State $b \gets r$
    \EndWhile
    \State \Return $a$
\end{algorithmic}
\end{algorithm}
\algcaption[alg:euclidean]{الگوریتم اقلیدس}
\end{latin}

\section{جدول مقایسه}

\begin{table}[h]
\centering
\caption{مقایسه پیچیدگی الگوریتم‌های مرتب‌سازی}
\label{tab:sorting-comparison}
\begin{tabular}{|c|c|c|c|c|}
\hline
\textbf{الگوریتم} & \textbf{بهترین} & \textbf{متوسط} & \textbf{بدترین} & \textbf{فضا} \\
\hline
حبابی & $O(n)$ & $O(n^2)$ & $O(n^2)$ & $O(1)$ \\
ادغامی & $O(n \log n)$ & $O(n \log n)$ & $O(n \log n)$ & $O(n)$ \\
سریع & $O(n \log n)$ & $O(n \log n)$ & $O(n^2)$ & $O(\log n)$ \\
هیپی & $O(n \log n)$ & $O(n \log n)$ & $O(n \log n)$ & $O(1)$ \\
\hline
\end{tabular}
\end{table}

%========================================
% PART 3: ADVANCED TOPICS
%========================================
\part{مفاهیم پیشرفته}

%========================================
% CHAPTER 5: Object-Oriented Programming
%========================================
\chapter{برنامه‌نویسی شیءگرا}
\label{chap:oop}

\section{کلاس و شیء}

برنامه‌نویسی شیءگرا\index{برنامه‌نویسی شیءگرا}
(\lr{OOP}) روشی برای سازمان‌دهی کد
بر اساس مفاهیم «کلاس»\index{کلاس} و «شیء»\index{شیء}
است.

\begin{definition}[کلاس]
    کلاس\index{کلاس!تعریف} قالبی برای ساخت اشیاء است
    که ویژگی‌ها (\lr{attributes}) و رفتارهای
    (\lr{methods}) مشترک را تعریف می‌کند.
    \label{def:class}
\end{definition}

\begin{latin}
\begin{minted}{python}
class BankAccount:
    """
    A simple bank account class.
    
    Attributes:
        account_number: Unique account identifier.
        owner: Name of the account owner.
        balance: Current account balance.
    """
    
    def __init__(
        self,
        account_number: str,
        owner: str,
        initial_balance: float = 0.0
    ):
        self.account_number = account_number
        self.owner = owner
        self.balance = initial_balance
    
    def deposit(self, amount: float) -> None:
        """Add money to the account."""
        if amount > 0:
            self.balance += amount
    
    def withdraw(self, amount: float) -> bool:
        """Withdraw money if sufficient balance."""
        if 0 < amount <= self.balance:
            self.balance -= amount
            return True
        return False


# Create and use accounts
account = BankAccount("IR-1234", "Ali", 1000)
account.deposit(500)
success = account.withdraw(200)
print(f"Balance: {account.balance}")  # 1300
\end{minted}
\end{latin}
\codecaption{کلاس حساب بانکی در پایتون}
\label{code:bank-account}

\section{مفاهیم اصلی OOP}

\begin{itemize}
    \item \textbf{کپسوله‌سازی}\index{کپسوله‌سازی}:
    پنهان‌سازی جزئیات پیاده‌سازی
    \item \textbf{وراثت}\index{وراثت}:
    ایجاد کلاس جدید بر اساس کلاس موجود
    \item \textbf{چندریختی}\index{چندریختی}:
    استفاده از یک رابط برای انواع مختلف
\end{itemize}

%========================================
% CHAPTER 6: Data Structures
%========================================
\chapter{ساختمان داده‌های پیشرفته}
\label{chap:data-structures}

\section{درخت جستجوی دودویی}

درخت جستجوی دودویی\index{درخت!جستجوی دودویی}
(\lr{BST}) ساختمان داده‌ای است که در آن
هر گره حداکثر دو فرزند دارد و مقدار
فرزند چپ کوچکتر و فرزند راست بزرگتر
از گره والد است.

\begin{figure}[h]
\centering
\begin{latin}
\begin{tikzpicture}[
  level distance=2cm,
  level 1/.style={sibling distance=4cm},
  level 2/.style={sibling distance=2cm}
]
  \node[treenode] {8}
    child {node[treenode] {3}
      child {node[treenode] {1}}
      child {node[treenode] {6}
        child {node[treenode] {4}}
        child {node[treenode] {7}}
      }
    }
    child {node[treenode] {10}
      child[missing]
      child {node[treenode] {14}
        child {node[treenode] {13}}
        child[missing]
      }
    };
\end{tikzpicture}
\end{latin}
\caption{درخت جستجوی دودویی}
\label{fig:bst}
\end{figure}

\section{صف و پشته}

\begin{table}[h]
\centering
\caption{مقایسه صف و پشته}
\label{tab:queue-stack}
\begin{tabular}{|c|C{4cm}|C{4cm}|}
\hline
\textbf{ویژگی} & \textbf{صف (\lr{Queue})} & \textbf{پشته (\lr{Stack})} \\
\hline
اصل & \lr{FIFO}\index{FIFO} & \lr{LIFO}\index{LIFO} \\
درج & انتها (\lr{enqueue}) & بالا (\lr{push}) \\
حذف & ابتدا (\lr{dequeue}) & بالا (\lr{pop}) \\
کاربرد & صف چاپ، \lr{BFS} & بازگشت، \lr{DFS} \\
\hline
\end{tabular}
\end{table}

%========================================
% CHAPTER 7: Graph Algorithms
%========================================
\chapter{الگوریتم‌های گراف}
\label{chap:graphs}

\section{جستجوی اول سطح}

\begin{latin}
\begin{algorithm}
\begin{algorithmic}[1]
    \Require Graph $G = (V, E)$, start vertex $s$
    \Ensure Shortest distances from $s$ to all vertices
    
    \ForAll{$v \in V$}
        \State $dist[v] \gets \infty$
    \EndFor
    \State $dist[s] \gets 0$
    \State Queue $Q \gets [s]$
    
    \While{$Q$ is not empty}
        \State $u \gets Q.\text{dequeue}()$
        \ForAll{neighbors $v$ of $u$}
            \If{$dist[v] = \infty$}
                \State $dist[v] \gets dist[u] + 1$
                \State $Q.\text{enqueue}(v)$
            \EndIf
        \EndFor
    \EndWhile
    \State \Return $dist$
\end{algorithmic}
\end{algorithm}
\algcaption[alg:bfs]{جستجوی اول سطح (\lr{BFS})}
\end{latin}

\section{الگوریتم دایکسترا}

\begin{latin}
\begin{algorithm}
\begin{algorithmic}[1]
    \Require Weighted graph $G = (V, E)$, source $s$
    \Ensure Shortest distances from $s$
    
    \ForAll{$v \in V$}
        \State $dist[v] \gets \infty$
    \EndFor
    \State $dist[s] \gets 0$
    \State $PQ \gets$ priority queue of $(dist[v], v)$
    
    \While{$PQ$ is not empty}
        \State $(d, u) \gets PQ.\text{extract\_min}()$
        \If{$d > dist[u]$} \textbf{continue} \EndIf
        \ForAll{neighbors $v$ of $u$}
            \If{$dist[u] + w(u,v) < dist[v]$}
                \State $dist[v] \gets dist[u] + w(u,v)$
                \State $PQ.\text{insert}(dist[v], v)$
            \EndIf
        \EndFor
    \EndWhile
    \State \Return $dist$
\end{algorithmic}
\end{algorithm}
\algcaption[alg:dijkstra]{الگوریتم دایکسترا}
\end{latin}

%========================================
% CHAPTER 8: Exercises
%========================================
\chapter{تمرین‌های پایانی}
\label{chap:exercises}

\section{تمرین‌های برنامه‌نویسی}

\begin{exercise}
    تابعی به نام \inlcode{is\_prime} بنویسید
    که یک عدد صحیح دریافت کند و مشخص کند
    آیا آن عدد اول است یا خیر.
    \label{exc:is-prime}
\end{exercise}

\begin{solution}
    \begin{latin}
\begin{minted}{python}
def is_prime(n: int) -> bool:
    """Check if n is a prime number."""
    if n < 2:
        return False
    for i in range(2, int(n**0.5) + 1):
        if n % i == 0:
            return False
    return True

# Test
print(is_prime(17))  # True
print(is_prime(100)) # False
\end{minted}
    \end{latin}
\end{solution}

\begin{exercise}
    الگوریتم جستجوی دودویی (\algref{alg:binary})
    را به صورت تابع پایتون پیاده‌سازی کنید.
    \label{exc:binary-search}
\end{exercise}

\begin{solution}
    \begin{latin}
\begin{minted}{python}
def binary_search(arr: list, target: int) -> int:
    """Binary search in sorted array."""
    left, right = 0, len(arr) - 1
    while left <= right:
        mid = left + (right - left) // 2
        if arr[mid] == target:
            return mid
        elif arr[mid] < target:
            left = mid + 1
        else:
            right = mid - 1
    return -1

# Test
arr = [1, 3, 5, 7, 9, 11, 13]
print(binary_search(arr, 7))  # 3
\end{minted}
    \end{latin}
\end{solution}

\section{تمرین‌های ریاضی}

\begin{exercise}
    ثابت کنید $\sqrt{2}$ عددی گنگ است.
    \label{exc:sqrt2}
\end{exercise}

\begin{proof}
    از برهان خلف استفاده می‌کنیم.
    فرض کنید $\sqrt{2}$ گویا باشد:
    $\sqrt{2} = \frac{a}{b}$ که $\gcd(a,b) = 1$.

    $2 = \frac{a^{2}}{b^{2}} \implies a^{2} = 2b^{2}$

    بنابراین $a^{2}$ زوج است $\implies$ $a$ زوج است.
    $a = 2k \implies a^{2} = 4k^{2} = 2b^{2}
    \implies b^{2} = 2k^{2}$

    بنابراین $b$ نیز زوج است.
    اما $\gcd(a,b) \geq 2$ با فرض $\gcd(a,b) = 1$
    تناقض دارد. $\therefore$ $\sqrt{2}$ گنگ است.
\end{proof}

%========================================
% CONCLUSION
%========================================
\chapter{نتیجه‌گیری}
\label{chap:conclusion}

در این کتاب، سفری از مبانی برنامه‌نویسی
تا الگوریتم‌های پیشرفته داشتیم.
از پایتون\index{پایتون} شروع کردیم،
با الگوریتم‌های جستجو\index{الگوریتم!جستجو}
و مرتب‌سازی\index{الگوریتم!مرتب‌سازی} ادامه دادیم،
و در نهایت به ساختمان داده‌های پیشرفته\index{ساختمان داده}
و الگوریتم‌های گراف\index{الگوریتم!گراف} رسیدیم.

برای مطالعه بیشتر، منابع معرفی شده
در بخش کتابنامه (\cite{knuth1984,cormen2009})
توصیه می‌شوند.

\begin{flushleft}
    \textbf{پایان}\\
    تابستان ۱۴۰۵
\end{flushleft}

%========================================
% BIBLIOGRAPHY
%========================================
\begin{latin}
\printbibliography[title={منابع و مراجع}]
\end{latin}

%========================================
% INDEX
%========================================
\printindex

\end{document}